\chapter[Encoder and convergence: more proofs]{Phases 7--8: encoder and convergence, the remaining proofs}\label{ch:cIV6}
\chapterbrief
{the notation of Chapters~\ref{ch:cIII3} and~\ref{ch:cIII4}.}
{to give the proofs of results of the encoder and convergence packages that were only listed in Chapter~\ref{ch:cIV2}: the structure of differential intertwiners, the interaction term with a potential, the spectral criterion for non-local laws, fractional and delay equations, the reduction to irreducible factors, the ambiguity in dimension $1+d$, and boundary conditions.}
{Dim.~B88--B93. They prove the corresponding statements of Chapters~\ref{PAP-ch:encoder} and~\ref{PAP-ch:convergence} of the paper.}

\section{The encoder}

\begin{dimcard}{B88}{ENC-K1p}{Differential intertwiners are Darboux chains}{\proved}{}
Let $J$ be an open interval, $q\in C^\infty(J;\C)$, $c\in\C$ and $W=\sum_{j=0}^nw_j\partial^j$ ($n\ge1$, $w_n\not\equiv0$) with $WL_q=(L_0+c)W$. Then (i) $w_n$ is constant; (ii) $N=\ker W$ has dimension $n$ and is $L_q$-invariant; (iii) there is a closed countable set $Z\subset J$ such that on each component of $J\setminus Z$, $W=w_nA_n\cdots A_1$ with $A_k=\partial-\psi_k'/\psi_k$, where $\psi_k$ is an eigenfunction of $L_{q_{k-1}}$ with eigenvalue $\lambda_k$, $q_0=q$, $q_k=q_{k-1}-2(\psi_k'/\psi_k)'$ and $q_n\equiv c$. Conversely every such Darboux chain gives a differential intertwiner (Crum). So, locally, a differential $W$ with $WL_q=(L_0+c)W$ exists iff $q$ is obtained from $c$ by a finite Darboux chain.
\end{dimcard}

\begin{proofidea}
Factor $W$ one first-order Darboux step at a time: an eigenvector of $L_q$ on the finite-dimensional kernel of $W$ gives the first factor, and the quotient is again an intertwiner, of order $n-1$ and for the transformed potential.
\end{proofidea}

\begin{proof}
\step{1}{(i) and (ii)}
The coefficient of $\partial^{n+1}$ in $WL_q-(L_0+c)W$ is $2w_n'$, so $w_n$ is a non-zero constant, hence $\dim\ker W=n$ on $J$. If $Wf=0$ then $WL_qf=(L_0+c)Wf=0$: $N$ is $L_q$-invariant.

\step{2}{(iii) One Darboux step}
$\C$ is algebraically closed, so $L_q|_N$ has an eigenvector $\psi_1$ with eigenvalue $\lambda_1$. On $J_1=J\setminus\psi_1^{-1}(0)$ (the zeros of a non-zero solution of a second-order linear ODE are isolated) put $g=\psi_1'/\psi_1$, $A_1=\partial-g$, $A_1^\sharp=-\partial-g$. Since $g'+g^2=\psi_1''/\psi_1=q-\lambda_1$:
\[
L_q-\lambda_1=A_1^\sharp A_1,\qquad A_1A_1^\sharp=L_{q_1}-\lambda_1\ \ (q_1=q-2g'),\qquad\text{hence}\qquad A_1L_q=L_{q_1}A_1 .
\]
$W\psi_1=0$, so right Euclidean division by the monic $A_1$ gives $W=W_1A_1+r$ with $r$ a function and $r\psi_1=0$, so $r=0$. From $W_1A_1L_q=(L_0+c)W_1A_1$ and $A_1L_q=L_{q_1}A_1$: $(W_1L_{q_1}-(L_0+c)W_1)A_1=0$. $A_1$ is onto $C^\infty$ of every subinterval (a first-order ODE), so $W_1L_{q_1}=(L_0+c)W_1$ with $W_1$ of order $n-1$.

\step{3}{Induction and the exceptional set}
Iterate: at $n=1$, $W_0=w_0$ is a function with $w_0L_{q_1}=(L_0+c)w_0$; the coefficient of $\partial$ gives $w_0$ constant and that of $\partial^0$ gives $q_1=c$. Let $Z=Z_1\cup\dots\cup Z_n$ with $Z_k$ the zeros of $\psi_k$ in $J\setminus(Z_1\cup\dots\cup Z_{k-1})$: each $Z_k$ is discrete in its open set, so $Z$ is countable, and closed because $Z_1\cup\dots\cup Z_k$ contains the closure of $Z_k$. (It is not claimed that $Z$ is discrete.)

\whyworks{a differential intertwiner has a finite-dimensional invariant kernel, which contains an eigenvector; dividing by the first-order operator built from it lowers the order by one.}
\end{proof}
\uses{B79.}

\begin{caution}[title=Priority]
The factorisation of an intertwiner into first-order Darboux steps is the subject of Andrianov--Sokolov (2003) and Sokolov (arXiv:0903.2835); the statement is almost certainly known. The global explicit classification for rational $V\to0$ is Duistermaat--Gr\"unbaum (1986) and holds only with decay at infinity (ENC-K1, from a cited source).
\end{caution}

\begin{dimcard}{B89}{ENC-A9}{The interaction term for Schr\"odinger with a potential}{\proved}{}
Let $T_1=-iH_0$ (kinetic), $T_2=-iV$, $U(t)=e^{-it(H_0+V)}$, $E^{\rm mult}(t)=U(t)-U_0(t)e^{-itV}$ and $V\in W^{2,\infty}(\R)$ real.
\begin{enumerate}[label=(\alph*)]
\item For $f\in H^1$, $t\ge0$: $\|E^{\rm mult}(t)f\|\le\frac{t^2}2(\|V''\|_\infty\|f\|+2\|V'\|_\infty\|f'\|)+\frac{2t^3}3\|V'\|_\infty^2\|f\|$; for $f\in H^2$, $E^{\rm mult}(t)f=\frac{t^2}2[H_0,V]f+o(t^2)$ with $[H_0,V]=-V''-2V'\partial$; and $E^{\rm mult}\equiv0$ near $t=0$ iff $V$ is constant.
\item ($V\equiv c$) $E^{\rm mult}=0$ and $E^{\rm add}=(e^{-ict}-1)(U_0-I)$.
\item ($V=Fx$ on $L^2(\R)$) $U(t)=e^{-iF^2t^3/12}e^{-itV/2}U_0(t)e^{-itV/2}$ (Strang, exact up to a phase), and $\|E^{\rm mult}(t)f\|\le|F|t^2\|f'\|+\frac23F^2t^3\|f\|$, $\|E^{\rm mult}(t)\|_{L^2\to L^2}=2$.
\end{enumerate}
The constant $2/3$ is optimal. (The first version wrote $F^2t^3/6$: false, because $e^{-itFx/2}$ translates $p$ by $Ft/2$; for $F=t=1$ and a Gaussian of width $25$ the true value is $0.6549$ against $0.1950$.)
\end{dimcard}

\begin{proof}
\step{1}{(a) Two Duhamel identities}
For $f\in H^2$ put $\mathcal S(s)=U_0(s)e^{-isV}$. Then $s\mapsto U(t-s)\mathcal S(s)f$ is $C^1$ with derivative $iU(t-s)[V,U_0(s)]e^{-isV}f$ (the two terms with $H_0$ cancel because $U_0$ commutes with $H_0$), so
\[
E^{\rm mult}(t)f=-i\int_0^tU(t-s)[V,U_0(s)]e^{-isV}f\dd s,\qquad [V,U_0(s)]g=i\int_0^sU_0(r)[H_0,V]U_0(s-r)g\dd r .
\]

\step{2}{The estimate and the leading term}
$\|[H_0,V]h\|\le\|V''\|\|h\|+2\|V'\|\|h'\|$, $U_0$ is unitary and commutes with $\partial$, and $\|(e^{-isV}f)'\|\le\|f'\|+s\|V'\|\|f\|$; hence
\[
\|E^{\rm mult}(t)f\|\le\int_0^ts\big[\|V''\|\|f\|+2\|V'\|(\|f'\|+s\|V'\|\|f\|)\big]\dd s,
\]
the stated bound; extend to $H^1$ by density. The double integral form shows $E^{\rm mult}(t)f/t^2\to\frac12[H_0,V]f$. If $E^{\rm mult}\equiv0$ then $[H_0,V]f=0$ for all $f\in H^2$; $f\equiv1$ on an interval gives $V''=0$, then $f'\ne0$ gives $V'=0$. Conversely for constant $V$, $e^{-itV}$ is a scalar.

\step{3}{(c) The Stark case in Fourier variables}
With $\hat f(p)=\int fe^{-ipx}$, $x=i\partial_p$ and the equation becomes $\partial_t\hat u-F\partial_p\hat u=-ip^2\hat u$; characteristics $p(s)=p+F(t-s)$ give
\[
\hat u(t,p)=\hat u_0(p+Ft)\exp\!\big(-i(p^2t+pFt^2+F^2t^3/3)\big).
\]
The Strang product gives $e^{-i(p^2t+pFt^2+F^2t^3/4)}\hat f(p+Ft)$ (since $e^{-itFx/2}$ shifts $p\mapsto p+Ft/2$), so the phase $F^2t^3/12$ closes the identity. $\|E^{\rm mult}(t)f\|$ is the norm of $(1-e^{i\theta(p)})$ applied after a translation, with $\theta$ affine and non-constant, whence the bound and the sup norm $2$.

\whyworks{the interaction of kinetic and potential terms is governed by the single commutator $[H_0,V]$, which is a first-order operator with coefficient $V'$: it vanishes iff $V$ is constant.}
\end{proof}
\uses{none.}

\begin{caution}[title=A warning on $E^{\rm add}$]
$E^{\rm add}$ does \emph{not} vanish for constant $V$: $E^{\rm add}=(e^{-ict}-1)(U_0(t)-I)$, with expansion $-\frac{t^2}2\{H_0,V\}+O(t^3)$. The quantity that measures the genuine interaction between kinetic and potential terms is $E^{\rm mult}$, and the right criterion is ``$E^{\rm mult}\equiv0\iff V$ constant''.
\end{caution}

\section{Locality, fractional laws and delays}

\begin{dimcard}{B90}{CONV-T6}{A spectral criterion for the annihilator of a non-local datum}{\proved}{}
Let $u(t,x)=\int_\R e^{a(\kappa)t+i\kappa x}\hat f(\kappa)\dd\kappa$ with $\hat f\in L^1$ of compact support, $a$ continuous on $\mathrm{supp}\,\hat f$, $E=\{\hat f\ne0\}$, $\Gamma_E$ the graph $\{(a(\kappa),i\kappa):\kappa\in E\}$.
\begin{enumerate}[label=(\arabic*)]
\item $\Ann(u)=I(\Gamma_E)=\{Q:Q(a(\kappa),i\kappa)=0\ \text{a.e. on }E\}$.
\item If $J\subseteq E$ is an open interval on which $a$ is real-analytic, $\Ann(u)\subseteq I_J=\{Q:Q(a,i\kappa)=0\ \text{on }J\}$, a prime ideal; $I_J=0$ if $a|_J$ is transcendental over $\C(\kappa)$, otherwise $I_J=(m_J)$ with $m_J$ irreducible. If $E$ is a finite union of such intervals, $\Ann(u)=\bigcap I_{J_i}$: zero if some branch is transcendental, otherwise principal, generated by the product of the distinct $m_{J_i}$.
\item Any multiplier $b=a$ a.e.\ on $E$ gives the same $u$: the non-local law is determined by $u$ only on $\Gamma_E$; if $a|_E$ is algebraic, the local law $m(\partial)u=0$ and the non-local law are indistinguishable.
\item If $u=\iint e^{i(\omega t+\kappa x)}\hat g$ with $\hat g\in L^1(\R^2)$ of compact support and $\hat g\ne0$ a.e.\ on a non-empty open set, then no real-analytic $m\not\equiv0$ satisfies $m(D)u=0$ on $\Omega$; in particular $\Ann(u)=0$.
\end{enumerate}
\end{dimcard}

\begin{proof}
\step{1}{(1) Differentiate under the integral}
$Q(\partial)u=\int Q(a(\kappa),i\kappa)e^{a(\kappa)t+i\kappa x}\hat f\dd\kappa$. For fixed $t$ this is, in $x$, the inverse Fourier transform of a compactly supported $L^1$ function, hence entire; if it vanishes on $\Omega$ it vanishes on a segment, hence on $\R$, and by injectivity of the Fourier transform $Q(a,i\kappa)e^{at}\hat f=0$ a.e. The converse is immediate.

\step{2}{(2) Primality and principality}
For $Q\in\Ann(u)$, $\kappa\mapsto Q(a(\kappa),i\kappa)$ is real-analytic on $J$ and zero a.e., hence identically zero. $I_J$ is the kernel of the homomorphism $R\to\mathcal O(J)$, $Q\mapsto Q(a,i\kappa)$ into an integral domain, so it is prime. If $a|_J$ is transcendental, $I_J=0$. Otherwise $I_J\ne0$ and not maximal (a maximal ideal contains $\xi-\xi_0$, impossible since $i\kappa$ is not constant on $J$), and a non-zero non-maximal prime of $\C[\tau,\xi]$ has height one, hence is generated by an irreducible polynomial (unique factorisation). For a finite union, $\Ann(u)=\bigcap I_{J_i}$ is the intersection of principal primes: generated by the product of the distinct generators.

\step{3}{(3) and (4)}
(3) is immediate from (1). (4) $m(D)u=\iint m\hat ge^{i(\omega t+\kappa x)}$ is entire; zero on $\Omega$ implies zero everywhere, so $m\hat g=0$ a.e., so $m=0$ a.e.\ on $U$, so $m\equiv0$ on $U$ and, by analyticity on $\R^2$, everywhere.

\whyworks{a datum remembers only the Zariski closure of the graph of its symbol; a law is local exactly when that graph is algebraic.}
\end{proof}
\uses{A2, A41.}

\begin{example}{Poisson, Klein--Gordon, telegraph, fractional}
Poisson $e^{-t|D|}$ ($a=-|\kappa|$, two intervals): $m=\tau-i\xi$ on $\kappa>0$ and $\tau+i\xi$ on $\kappa<0$; the product $\tau^2+\xi^2$: $u$ is harmonic. $e^{-t\sqrt{1-\partial^2}}$ gives $\tau^2+\xi^2-1$ (elliptic Klein--Gordon). Exponential memory $u_t=\int_{-\infty}^te^{-a(t-s)}u_{xx}(s)\dd s$ gives $\tau(\tau+a)=\xi^2$ (telegraph). A fractional $a=-\kappa^{p/q}$ on an interval of $(0,\infty)$ with $\gcd(p,q)=1$ has minimal polynomial dividing $(-\tau)^q-(-i\xi)^p$.
\end{example}

\begin{dimcard}{B91}{CONV-P6, CONV-P7}{Irrational fractional power and delay: $\Ann=0$; locality is a property of the pair $(u,\mathcal K)$}{\proved}{}
Let $J\subset(0,\infty)$ be an open interval and $\hat f\ne0$ a.e.\ on $J$.
\begin{enumerate}[label=(\alph*)]
\item If $a(\kappa)=-\kappa^\alpha$ with $\alpha>0$ irrational, $\Ann(u)=0$.
\item If $\sigma>0$, $J\subset(0,(\sigma e)^{-1/2})$ and $a(\kappa)=W_0(-\sigma\kappa^2)/\sigma$ (principal Lambert branch), $u$ solves $u_t(t,x)=u_{xx}(t-\sigma,x)$ and $\Ann(u)=0$.
\item (the delay is local after augmentation) $w(t,\theta,x)=u(t-\theta,x)$ solves $w_t+w_\theta=0$ with boundary coupling $\partial_tw(t,0,x)=\partial_x^2w(t,\sigma,x)$; conversely $w=w(t-\theta,0,x)$ and $u=w(\cdot,0,\cdot)$ solves the delay equation.
\item (P7) Without the non-triviality condition every $u$ lies in $\mathcal L^+_{\mathcal K}$ (take $Y=\R$, $S=\{\partial_yv=0\}$, $\gamma(v)=v(\cdot,\cdot,0)$). With it: the delay is in $\mathcal L^+_{\rm transport}\setminus\mathcal L$; exponential memory is in $\mathcal L\cap\mathcal L^+_{\rm ODE}$; a datum with a two-dimensional spectrum is in $\mathcal N_{\mathcal K}$ for $\mathcal K$ real-analytic translation-invariant multipliers.
\end{enumerate}
\end{dimcard}

\begin{proof}
\step{1}{(a) Independence of exponentials}
If $Q=\sum c_{ij}\tau^i\xi^j\ne0$ vanishes on the graph, $H(\kappa)=\sum c_{ij}(-1)^ii^j\kappa^{i\alpha+j}$ is analytic on $(0,\infty)$ and zero on $J$, so zero on $(0,\infty)$. The exponents $i\alpha+j$ are distinct for distinct $(i,j)$ ($\alpha$ irrational); with $\kappa=e^s$ the functions $e^{(i\alpha+j)s}$ are exponentials with distinct frequencies, hence independent (A2): all $c_{ij}=0$. Apply B90(2).

\step{2}{(b) Lambert}
For $y\in[-1/e,0)$, $W_0(y)e^{W_0(y)}=y$ is real, real-analytic and non-constant on $(-1/e,0)$. With $a=W_0(-\sigma\kappa^2)/\sigma$: $a=-\kappa^2e^{-\sigma a}$, so $u_t=\int ae^{at+i\kappa x}\hat f$ equals $u_{xx}(t-\sigma)=\int(-\kappa^2)e^{a(t-\sigma)+i\kappa x}\hat f$. For transcendence: if $Q(a,i\kappa)=0$ then $a$ is algebraic over $\C(\kappa)$, $a\ne0$ on $J$, so $e^{\sigma a}=-\kappa^2/a$ is algebraic over $\C(\kappa)$; as $a$ is non-constant, $\kappa$ is algebraic over $\C(a)$, so $e^{\sigma a}$ would be algebraic over $\C(a)$, impossible ($a$ transcendental over $\C$ and $e^{\sigma a}$ transcendental over $\C(a)$). So $I_J=0$ and $\Ann(u)=0$ by B90(2).

\step{3}{(c) and (P7)}
(c) $w_t+w_\theta=u_t(t-\theta)-u_t(t-\theta)=0$; $w(t,0,x)=u$, $w(t,\sigma,x)=u(t-\sigma,x)$ gives the coupling; conversely $w$ is constant along $(t+s,\theta+s)$. (P7)(a) is the given trivial system; (b): the traces of solutions of the delay system are only the solutions of the delay equation (for example $e^t\sin x$ is not one: $1\ne-e^{-\sigma}$), and $\Ann(u)=0$ by (b); memory: $v=\int_{-\infty}^te^{-a(t-s)}u_{xx}(s)\dd s$ solves $v_t=-av+u_{xx}$ and $u_t=v$; the two-dimensional spectrum case is B90(4).

\whyworks{``local or non-local'' is not a property of $u$ but of the pair ($u$, a declared class of laws): without a non-triviality condition the class of augmented laws is empty of content.}
\end{proof}
\uses{A2, B90.}

\section{Dimension $1+d$}

\begin{dimcard}{B92}{CONV-P8, CONV-D1}{Reduction to irreducible factors, and the ambiguity in $1+d$}{\proved}{}
\textbf{(P8)} Let $P=\prod_iP_i\in R_n$ be reduced ($P_i$ pairwise non-associate irreducibles), $F$ any $R_n$-module, $u\in F$ with $Pu=0$, $w_i=(P/P_i)u$. Then $\Ann(u)=(P)\iff\Ann(w_i)=(P_i)$ for every $i$.

\textbf{(D1)} Let $\Omega\subset\R^n$ be any open set and $\mathcal A=C^\infty(\Omega)$ or $\mathcal D'(\Omega)$. (a) For $Q_1,Q_2\in R_n$, each $\zeta\in V(Q_1,Q_2)$ gives $e_\zeta=e^{\langle\zeta,z\rangle}\in\ker Q_1\cap\ker Q_2$, with distinct $\zeta$ independent; for $n\ge3$, $\dim R_n/(Q_1,Q_2)<\infty\iff(Q_1,Q_2)=R_n\iff V(Q_1,Q_2)=\emptyset$, otherwise $\ker Q_1\cap\ker Q_2$ is infinite-dimensional. (b) The ambiguity space $\mathcal K=\{(v_j)\in\bigoplus\ker Q_j:\sum v_j=0\}$ is $0$ iff the $Q_j$ are pairwise co-maximal; for $n\ge3$ and $\mathcal K\ne0$, $\dim\mathcal K=\infty$. (c) (a non-EP exceptional class) in $\R^3$ with $P_i=\tau-\xi_1$, $w=f(t+x_1)$, $f$ not exponential-polynomial, has $\Ann(w)=(\tau-\xi_1,\xi_2)\supsetneq(P_i)$ and $\dim R_3w=\infty$.

So in $1+d$, $d\ge2$, $\delta(P)\in\{0,\infty\}$, with $\delta=0$ iff the factors are pairwise co-maximal, on every open set without convexity.
\end{dimcard}

\begin{proof}
\step{1}{(P8)}
($\Rightarrow$) $P_iw_i=Pu=0$, so $(P_i)\subseteq\Ann(w_i)$. If $Qw_i=0$ then $Q(P/P_i)\in\Ann(u)=(P)$, so $P_i\mid Q(P/P_i)$; $P_i$ is prime and does not divide $P/P_i$, so $P_i\mid Q$. ($\Leftarrow$) $(P)\subseteq\Ann(u)$; if $Qu=0$ then $Qw_i=(P/P_i)Qu=0$, so $P_i\mid Q$ for all $i$, and distinct primes give $P\mid Q$.

\step{2}{(D1)(a)}
$Q(\partial)e_\zeta=Q(\zeta)e_\zeta$; independence is A2 restricted to a generic line on which the linear forms $\langle\zeta,\cdot\rangle$ stay distinct. For $n\ge3$ and $V(Q_1,Q_2)\ne\emptyset$ every irreducible component has dimension $\ge n-2\ge1$ (Krull's principal ideal theorem, twice), so $V$ is infinite, $R_n/(Q_1,Q_2)$ is infinite-dimensional (Nullstellensatz) and there are infinitely many independent exponentials in the intersection. If $V=\emptyset$, $(Q_1,Q_2)=R_n$.

\step{3}{(D1)(b) and (c)}
(b) If the $Q_j$ are pairwise co-maximal then $Q_i$ and $\prod_{j\ne i}Q_j$ are co-maximal: $1=AQ_i+B\prod_{j\ne i}Q_j$. From $\sum v_j=0$, applying $\prod_{j\ne i}Q_j$ gives $\prod_{j\ne i}Q_jv_i=0$, and with $Q_iv_i=0$: $v_i=AQ_iv_i+B\prod_{j\ne i}Q_jv_i=0$. Conversely if $(Q_i,Q_j)\ne R_n$ pick $\zeta\in V(Q_i,Q_j)$: $(v_i,v_j)=(e_\zeta,-e_\zeta)$ is a non-zero element of $\mathcal K$, and for $n\ge3$ infinitely many independent ones exist. (c) $(\tau-\xi_1)w=f'-f'=0$, $\xi_2w=0$; for $Q\in\Ann(w)$ reduce modulo $(\tau-\xi_1,\xi_2)$ to $q(\xi_1)$, and $q(\partial)f=0$ forces $q=0$ because $f$ is not exponential-polynomial; $R_3w\cong\C[\xi_1]$.

\whyworks{in more than two variables two coprime polynomials still meet along a variety of positive dimension, so the overlap of two factors is either empty or infinite.}
\end{proof}
\uses{A2, A13.}

\section{Boundary conditions}

\begin{dimcard}{B93}{CONV-B1, CONV-B2}{Boundary conditions are invisible to $\Ann(u)$; stability with a dictionary of traces}{\proved}{}
\textbf{(B1)} For $u$ real-analytic on connected $\Omega$, $\Ann(u)$ is determined by the germ of $u$ at any interior point (A41(a)). On $(0,L)$ with $\mu=\pi/L$, $\kappa_0=D\mu^2$, $u_D=e^{-\kappa_0t}\sin\mu x$ (Dirichlet) and $u_N=e^{-\kappa_0t}\cos\mu x$ (Neumann) have the same $\Ann=(\tau+\kappa_0,\ \xi^2+\mu^2)$.

\textbf{(B2)} Let $\Omega$ be bounded Lipschitz, $\Gamma\subseteq\partial\Omega$ relatively open, $\gamma_\Gamma:H^1(\Omega)\to L^2(\Gamma)$ the trace (constant $C_{\rm tr}$), $\Phi_\Gamma(u)\lambda=\gamma_\Gamma(\sum_{|\alpha|\le n_0}\lambda_\alpha\partial^\alpha u)$, with $\Phi_\Gamma(u)\lambda^*=0$, $|\lambda^*|=1$, $\sigma_2(\Phi_\Gamma(u))\ge\rho_\Gamma>0$. If $\|\tilde u-u\|_{H^{n_0+1}}\le\delta$ and $\tilde\lambda$ minimises $\|\Phi_\Gamma(\tilde u)\lambda\|$ on $|\lambda|=1$, then $\sin\angle(\tilde\lambda,\lambda^*)\le2C_{\rm tr}\sqrt{D+1}\,\delta/\rho_\Gamma$; the misfit certificate and discrepancy rule of A43 carry over with $\delta$ replaced by $C_{\rm tr}\sqrt{D+1}\,\delta$.
\end{dimcard}

\begin{proof}
\step{1}{(B1)}
Both $\tau+\kappa_0$ and $\xi^2+\mu^2$ kill both functions; for each $\dim R\cdot u=2$ ($u$ and $u_x$ are independent: their Wronskian in $x$ is non-zero), and $\dim R/(\tau+\kappa_0,\xi^2+\mu^2)=2$. Since $(\tau+\kappa_0,\xi^2+\mu^2)\subseteq\Ann(u)$ and $R\cdot u\cong R/\Ann(u)$, the inclusion is an equality.

\step{2}{(B2)}
$\|(\Phi_\Gamma(\tilde u)-\Phi_\Gamma(u))\lambda\|\le\sum|\lambda_\alpha|C_{\rm tr}\|\partial^\alpha e\|_{H^1}\le C_{\rm tr}|\lambda|(\sum_\alpha\|\partial^\alpha e\|_{H^1}^2)^{1/2}$; each $\partial^\gamma$ with $|\gamma|\le n_0+1$ appears at most $D+1$ times (as $\partial^\alpha$ itself or as $\partial_j\partial^\alpha$), so the sum is $\le(D+1)\|e\|^2_{H^{n_0+1}}$. Then the proof of A43 (T4a) word for word.

\whyworks{the annihilator is a germ-level object, and a boundary condition selects among the solutions without changing the equation they satisfy.}
\end{proof}
\uses{A41, A43.}

\begin{dimcard}{B155}{CONV-B3}{Tangential obstruction: no boundary margin above the minimal level}{\pw}{ch:convergence}
Let $\Gamma=\{x_1=c\}\cap\bar\Omega$ be a flat boundary piece transversal to $t$, $u\in C^{n_0+2}$ up to $\Gamma$, and $\Phi_\Gamma(u)\lambda=\gamma_\Gamma(\sum_{|\alpha|\le n_0}\lambda_\alpha\partial^\alpha u)$. If $\lambda^*\in\ker\Phi_\Gamma(u)$ corresponds to $Q^*$ with $\deg Q^*\le n_0-1$, then $\tau Q^*$ is also in the kernel, $\dim\ker\Phi_\Gamma(u)\ge2$ and $\sigma_2(\Phi_\Gamma(u))=0$: no margin. With Dirichlet data and $n_0\ge1$ the kernel contains $1$ and $\tau$.
\end{dimcard}

\begin{proof}
\step{1}{The trace of $Q^*u$ vanishes}
$w=Q^*(\partial)u$ is $C^1$ up to $\Gamma$ and $w|_\Gamma=0$.

\step{2}{A tangential derivative preserves it}
$\partial_t$ is tangent to $\Gamma$, so $(\partial_tw)|_\Gamma=\partial_t(w|_\Gamma)=0$: $(\tau Q^*)(\partial)u$ has zero trace.

\step{3}{Independence}
$Q^*$ and $\tau Q^*$ have different degrees, hence are linearly independent in $R_{\le n_0}$. So $\dim\ker\ge2$ and $\sigma_2=0$.

\whyworks{a boundary condition is automatically accompanied by all its tangential consequences.}
\end{proof}
\uses{none.}

\begin{caution}[title={Reading}]
Identifying a boundary condition with a margin requires the minimal level (the rule of T5), or a dictionary of normal derivatives only, or the quotient by tangential consequences. Example: $u=e^{-t}\sin x$ on $x=0$ has traces $(0,0,e^{-t})$ of $(1,\partial_t,\partial_x)$.
\end{caution}

